\documentclass[11pt]{article}
\usepackage[margin=1in]{geometry}
\usepackage[T1]{fontenc}
\usepackage{lmodern,microtype,amsmath,amssymb,amsthm,graphicx,booktabs,tabularx}
\usepackage[colorlinks=true,linkcolor=blue,urlcolor=blue]{hyperref}
\setlength{\parindent}{0pt}
\setlength{\parskip}{4pt}
\newtheorem{proposition}{Proposition}
\newcommand{\dd}{\mathrm d}
\begin{document}
\begin{center}
{\Large Housing allocation, fertility, and population}\par
\smallskip
{\small Proposed results following the existing household setup}
\end{center}

\section*{Housing allocation}

Borrowing limits can leave housing with old owners even when young households
would give up more consumption to use it. In an equilibrium of the model,
write $x_i=c_i-\chi n_i$ and $s_i=h_i-\kappa n_i$ for
the goods and space available to young adults. The young household's and old
owner's marginal values of housing, in current consumption goods, are:
\begin{equation}
 MV_i^Y=\frac{\alpha x_i}{s_i},\qquad
 MV_j^O=\frac{\gamma c_j^2}{h_j^2}.
 \label{eq:core_values}
\end{equation}
These values measure what each household would pay for a little more space.

The planner reallocates housing and consumption, retaining the physical
limits on each tenure and the total housing stock. It may relax private
financing restrictions. Each household's fertility, tenure, estate, and
future real allocation remain fixed, and existing financial claims are honored.

\begin{proposition}[Housing misallocation]
Suppose positive masses of young owners can occupy more housing and old
owners can release some. If $MV_i^Y>MV_j^O$ on matched groups, a sufficiently
small reallocation toward the young can make them better off and leave
every other household equally well off, in the absence of real reallocation costs.
\end{proposition}

\textit{Proof.} Give the young household $\epsilon$ units of the old owner's
housing. Keeping the old owner's utility fixed requires extra consumption:
\begin{equation}
 D_j(\epsilon)=c_j^2\left[
 \left(\frac{h_j^2}{h_j^2-\epsilon}\right)^\gamma-1\right].
 \label{eq:core_compensation}
\end{equation}
Take this amount from the young household. Goods and housing balance, and
its utility gain has derivative $(MV_i^Y-MV_j^O)/x_i>0$ at zero.
A small change preserves feasibility.
Financial settlement replaces the seller's lost estate title with an equal
bond claim; the buyer's later sale of the extra title offsets its financing.
Future allocations, existing creditors, and aggregate external assets are
unchanged.\hfill$\square$

\begin{center}
\includegraphics[width=.70\textwidth]{output/model/simplified_olg_amendments/simple_allocation_figure.pdf}
\end{center}
{\small\textbf{Figure 1.} Housing moves toward the young while the old
owner is compensated. The shaded gap indicates the direction of improvement;
its area is not a welfare measure. Equality of values gives the best allocation
along this particular compensated comparison.}

\newpage
\section*{Housing access and fertility}

The borrowing limit can produce the inequality in Proposition 1. With
positive saving, slack physical owner limits, and interior old-age choices,
a strictly restrictive down payment gives $MV_i^Y>q r_t$. An old owner with
slack retention and estate bounds has $MV_j^O=q r_t$. Young buyers need cash
at purchase; old owners already hold their homes.

Rental segmentation limits another way to obtain space. A household may be
unable to rent a larger home and unable to fund its purchase. The ownership
taste in the model also matters: it can make constrained ownership preferable
even when larger rentals are available. Thus both access restrictions can
coexist with misallocation, although a binding rental cap is not necessary
for the owner-to-owner result. A renter already at a physical cap cannot
receive more housing in that tenure under the same planner restrictions.

Real reallocation costs have a simple implication. If moving $\epsilon$
housing uses $C(\epsilon)$ goods, with $C(0)=0$ and $C'(0)=K_{ijt}\geq0$, the condition
becomes $MV_i^Y>MV_j^O+K_{ijt}$. A positive fixed cost instead requires a finite
reallocation whose gain covers that cost.

\textbf{A fertility prediction.} More space makes children easier to
accommodate; paying for it reduces the goods available to the family.
Consider the young owner's conditional problem with positive saving and
interior old-age choices. Suppress household and date indices, and let
$k=\beta(1+\gamma+\omega_B)$. The fertility condition and the response of
the lifetime budget to more housing are:
\begin{equation}
 \frac{\vartheta}{n}=\frac{\chi}{x}+\frac{\alpha\kappa}{s},
 \qquad (1+k)\dd x=-\chi\dd n-p\dd h.
 \label{eq:core_fertility_system}
\end{equation}
Here $p$ is the net lifetime payment for a marginal unit of current housing,
including accompanying transfers. With prices and other resources fixed,
$p=q r_t$. Differentiating the first equation using the second gives:
\begin{equation}
 \frac{\dd n}{\dd h}>0
 \quad\Longleftrightarrow\quad
 \frac{\alpha\kappa}{s^2}>
 \frac{\chi p}{(1+k)x^2}.
 \label{eq:core_fertility_sign}
\end{equation}
Fertility rises when the space benefit exceeds the resource cost. This is a
household response while tenure and the binding constraints remain unchanged.
Saving and old-age choices are reoptimized. Proposition 1 instead holds
future real allocations fixed.

A sufficient condition can be written without prices or a borrowing
multiplier. At stationary prices, with zero property tax and a strictly
binding down-payment limit, it is:
\begin{equation}
 \frac{(1-q)b}{(1-\phi)(y+b)}
 \geq\frac{\alpha}{1+\beta(1+\gamma+\omega_B)+\alpha}.
 \label{eq:core_primitive}
\end{equation}
The left side is housing expenditure at the purchase limit relative to
lifetime resources. Under the stated household conditions, this restriction
ensures a fertility increase for any $p<MV^Y$. The full equilibrium response
also depends on prices, tenure choices, and subsequent cohort sizes.

\newpage
\section*{Fertility and population over time}

A temporary fertility increase can leave a permanently larger population.
Let $Y_t$ and $O_t$ count young and old households, $\bar n_t$ denote mean
young fertility, and $\nu$ convert children into entering households. Two
paths starting from the same young cohort satisfy:
\begin{equation}
 \frac{Y_t^1}{Y_t^0}=
 \prod_{s=0}^{t-1}\frac{\bar n_s^1}{\bar n_s^0},
 \qquad O_t^j=Y_{t-1}^j.
 \label{eq:core_population}
\end{equation}
Population reflects cumulative fertility differences. Both positive steady
states have replacement fertility $\bar n^*=1/\nu$. Housing clearing gives:
\begin{equation}
 N_{\rm hh}^*=2Y^*=
 \frac{2\bar H}{\bar h^{Y*}+\bar h^{O*}}.
 \label{eq:core_stationary_population}
\end{equation}
With a fixed stock, a larger population requires less housing per household
on average. This stationary ranking also applies to persons under unchanged
household and child-counting conventions.

The model admits local equilibrium transitions in which a small permanent
increase in the financed share raises initial fertility and final population.
The result is established near an economy with all owners, zero child goods
cost, and zero property tax. It permits heterogeneous entrant income and
wealth and small positive renter shares and costs, provided the same household
constraints continue to bind. The purchase limit remains restrictive, saving
is positive, and physical owner, old retention, and estate bounds are slack.
Convergence and the
initial fertility sign require separate parameter conditions.\footnote{The
conditions and proof are in Appendix B of the accompanying longer
assessment and sections 5--10 of its supporting proof. The permitted reform
size and positive-parameter neighborhood are local and have no numerical bound.}

\begin{center}
\includegraphics[width=\textwidth]{output/model/simplified_olg_amendments/local_transition_analytical_figure.pdf}
\end{center}
{\small\textbf{Figure 2.} Analytical first-order equilibrium response in
the illustrative all-owner limit, per unit increase in $d=b/(1-\phi)$.
Fertility rises most in the second young cohort and returns to replacement;
population converges to a higher level. All-date signs hold for small finite
reforms in this example; the nearby mixed-tenure result establishes initial
fertility and final population signs.}

In this limit, the stationary price rises in proportion to borrowing capacity.
Each young household's housing and fertility return to their original levels.
Older households use less housing, and the stock accommodates more households.
Stationary entrants of each type have lower lifetime utility. This transition
illustrates a demographic response to credit; the compensated welfare gain
in Proposition 1 is a separate comparison.

\newpage
\section*{For discussion: how much theory to retain}

\textbf{The allocation result can carry the main claim.} It identifies a
feasible welfare gain, explains why housing can remain with the old, and
requires one short proof. The planner's borrowing permission is the one
already accepted. A complete first-best allocation and a market-based
implementation remain separate questions.

\textbf{The transition proof is useful, but its supporting assumptions are
substantive.} The stability inequality itself is fairly simple. It does not
remove the following restrictions:

\renewcommand{\arraystretch}{1.2}
\begin{tabularx}{\textwidth}{@{}p{.29\textwidth}X@{}}
\toprule
Assumption & Economic consequence\\
\midrule
Near-all ownership and small child goods costs
& The proved neighborhood gives limited scope for tenure switching and the
goods cost of children. Neither its width nor a large renter share is covered.\\
Old owners can sell; estate bounds are slack
& The interior formula requires slack retention and, at zero property tax,
$(1-q)\omega_B>q\gamma$. A binding bound changes old housing demand.\\
Fixed entrant endowments and world bond price
& Larger cohorts receive the same distribution of income and wealth. Wages,
interest rates, and inherited wealth do not adjust to the larger population.\\
\bottomrule
\end{tabularx}

A small reform is a reasonable scope for an illustrative result. Uniformly
unchanged household constraints keep the local proof manageable. The stronger
limitations are the economic starting point and the fixed income and wealth
of new entrants. This is an existence result for a class of economies, with
explicit conditions; its assumptions have not been disciplined by data.

\textbf{The welfare loss is not peculiar to the plotted parameters.}
At the all-owner limit with $\chi=\tau^p=0$, let $d_i=b_i/(1-\phi)$ and
hold the entrant distribution fixed. For each type, the stationary owner value
satisfies:
\begin{equation}
 \frac{\dd W^{O*}(i)}{\dd\phi}
 =-\frac{(1-q)d_i/x_i^*+\beta\gamma}{1-\phi}<0.
 \label{eq:core_stationary_welfare}
\end{equation}
Higher prices use more of the household's resources without increasing its
stationary young housing. Old-age housing also falls. The formula holds on
the stated owner branch with fixed preferences; ownership tastes are held
fixed in the comparison. It compares the same type in two steady states,
without assigning a social value to the change in population or evaluating
the transitional cohorts.

\textbf{My recommendation.} Use the allocation result and the conditional
fertility condition as the core. Keep both figures available, presenting the
transition as a separate illustrative extension and putting its proof in the
appendix. Its value is the connection between temporary fertility and lasting
population, with the price response made explicit. Whether that extension
earns main-text space is the next useful discussion. Constrained efficiency
and public lending can remain later branches. No model convention or planner
power is changed by this proposed presentation.

\end{document}
